%&pdftex
\hsize=13.8true cm
\vsize=21.3true cm % \vsize=21.2true cm
\font\tit=cmssbx10 scaled 1200
\font\yo=cmcsc10
\font\cur=cmss10
\font\msym=msbm10
\def\R{\hbox{\msym R}}
\def\N{\hbox{\msym N}}
\def\'#1{\if#1i{\accent"13\i}\else{\accent"13 #1}\fi}
%\nopagenumbers

\bigskip


\centerline{\tit ACOTACI\'ON EN NORMA DE SERIES DE FOURIER-NEUMANN%
\footnote{}{\rm PUBLICADO EN:
{\it Actas XV Jornadas Luso-Espanholas de Matem\'atica}, Vol.~II
(\'Evo\-ra, Portugal, 1990), 119--124,
Universidade de \'Evora, \'Evora, 1991.}}

\bigskip

\centerline{\yo Juan L. Varona Malumbres}
\vskip 3pt
\centerline{\cur Colegio Universitario de La Rioja (Universidad de Zaragoza)}

\bigskip

{\narrower\noindent{$\underline {\rm Abstract}$.} If we take $\alpha > -1$,
the system $x^{-1/2}J_{\alpha+2n+1}(x)$, $n = 0, 1, 2, \dots$ is orthogonal on
$(0,\infty)$. We study the boundedness on $L^p(0,\infty)$
of the partial sum operator associated with the Fourier series, which is related
with its mean convergence. \smallskip}

\bigskip
\noindent{\sl Clasificaci\'on A.M.S. (1980): 42C10}
\bigskip

Si denotamos mediante $J_\mu(x)$ la funci\'on de Bessel de orden $\mu$ y tomamos
$\alpha > -1$ fijo, se verifica la relaci\'on de ortogonalidad [7]
$$\int_0^\infty  {J_{\alpha +2n+1}(x)J_{\alpha +2m+1}(x)x^{-1}\,dx} = {{\delta _{nm}}
\over {4n+2\alpha +2}},\quad n,m\in \N.$$

Por lo tanto, el sistema
$$j_n(x)=\sqrt {{{4n+2\alpha +2} \over x}}J_{\alpha+2n+1}(x), \quad n\in \N,$$
es ortonormal en $L^2 = L^2((0,\infty),dx)$.

Llamemos ahora $B_2$ a la clausura lineal en $L^2$ de
$\{j_n(x)\}_{n=0}^\infty$. La teor\'ia de espacios de Hilbert demuestra que,
para cada $f \in B_2$, su serie de Fourier
$$\sum\limits_{n=0}^\infty  {a_n(f)j_n(x)},\qquad a_n(f) = \int_0^\infty 
{j_n(x)f(x)\,dx}$$
converge a $f$ en la norma de $L^2$.

Del mismo modo, si llamamos $B_p$, $1<p<\infty$, a la clausura lineal de
$\{j_n(x)\}_{n=0}^\infty$ en $L^p = L^p((0,\infty),dx)$, nos preguntamos cu\'ando
la serie de Fourier de cada $f \in B_p$ converge a $f$ en la norma de $L^p$. Es
decir, si denotamos
$$S_n(f,x)=\sum\limits_{k=0}^n {a_k(f)j_k(x)},\qquad a_k(f) = \int_0^\infty 
{j_k(x)f(x)\,dx}  \eqno (1)$$
queremos saber cu\'ando $S_nf$ existe para cada $f \in B_p$ y $S_nf
\rightarrow f$ en $\Vert ~ \Vert_p$. Como los desarrollos del tipo $\sum_{m\ge 0}
{a_mJ_{\alpha +m}(x)}$ se llaman series de Neumann, a la expresi\'on (1) la
denominamos sumas parciales de la serie de Fourier-Neumann. (Un estudio de la
convergencia puntual de estas series puede verse en [8] y [9].)

Es bien conocido [6] que una condici\'on equivalente a la convergencia de $S_nf$
en $B_p$ es que se verifique la acotaci\'on uniforme $\Vert S_nf \Vert _p \le C
\Vert f \Vert _p,\ f \in B_p$, donde, de aqu\'i en adelante, $C$ denota una
constante independiente de $n$ pero que puede no ser la misma cada vez que
aparezca. Realmente, no nos preocuparemos de describir el espacio $B_p$ y
estudiaremos una condici\'on un poco m\'as exigente:
$$\Vert S_nf \Vert _p \le C \Vert f \Vert _p,\qquad f \in L^p. \eqno(2)$$

Es f\'acil demostrar por dualidad [5] que una condici\'on necesaria para la
acotaci\'on uniforme (2) es
$$\Vert j_n \Vert_p \Vert j_n \Vert_q \le C \eqno(3)$$
donde $1/p+1/q=1$.

En primer lugar, tomando $n=0$ en (3) es claro que obtenemos la
condici\'on necesaria
$$x^{-1/2}J_{\alpha+1}(x) \in L^p \cap L^q. \eqno(4)$$

Utilizando $J_{\mu}(x) = O(x^{-1/2}),\ x \rightarrow \infty,$ y la estimaci\'on
asint\'otica $J_{\mu}(x) \sim x^{\mu},\ x \rightarrow 0$, es inmediato comprobar
que
$${\rm cuando }\ \alpha \ge -1/2,\qquad (4) \Longleftrightarrow 1 < p <\infty,
\eqno(5)$$
$${\rm cuando }\ -1 < \alpha < -1/2,\qquad (4) \Longleftrightarrow {2 \over
{3+2\alpha }} < p < {2 \over {-1-2\alpha }}. \eqno(6)$$

Adem\'as, estimaciones asint\'oticas m\'as precisas para
$J_{\alpha+2n+1}(x)$ permiten demostrar (como en [1])
$$(3) \Longrightarrow 4/3 < p < 4. \eqno(7)$$

Una vez analizadas las condiciones necesarias para (2), vamos a dedicar
nuestros esfuerzos a demostrar que estas mismas condiciones sobre
$\alpha$ y sobre $p$ son suficientes para que se verifique (2).

Es claro que $$S_n(f,x) = \int_0^\infty  {K_n(x,y)f(y)\,dy},$$
donde
$$K_n(x,y) = \sum\limits_{k=0}^n
{j_k(x)j_k(y)} = x^{-1/2}y^{-1/2}\sum\limits_{k=0}^n {(4k+2\alpha +2)J_{\alpha
+2k+1}(x)J_{\alpha +2k+1}(y)}.$$

La expresi\'on anterior no nos permite estudiar directamente el
comportamiento de la serie de Fourier-Neumann, sino que necesitamos una
descomposici\'on del n\'ucleo $K_n(x,y)$ adecuada para nuestros
prop\'ositos:

\bigskip

\proclaim LEMA 1. Se verifica
$$K_n(x,y)={{x^{1/2}y^{1/2}} \over {x^2-y^2}}\left\{ {xJ_{\alpha +1}(x)J_\alpha
(y)-yJ_\alpha (x)J_{\alpha +1}(y)} \right\}$$
$$\null+ {{x^{1/2}y^{1/2}} \over {x^2-y^2}}\left\{ {xJ'_{\alpha +2n+2}(x)J_{\alpha
+2n+2}(y) - yJ_{\alpha +2n+2}(x)J'_{\alpha +2n+2}(y)} \right\}.$$

{\it Demostraci\'on}. En [8] se prueba que
$$\phi _\mu(x,y) = \sum\limits_{k=0}^\infty {(4k+2\mu +2)J_{\mu
+2k+1}(x)J_{\mu +2k+1}(y)}$$
cumple
$$\phi _\mu(x,y) = {{xy} \over {x^2-y^2}}\left\{ {xJ_{\mu +1}(x)J_\mu (y)-yJ_\mu
(x)J_{\mu +1}(y)} \right\}$$
y es claro que
$$K_n(x,y) = (xy)^{-1/2}\left\{ {\phi _\alpha (x,y)-\phi _{\alpha +2n+2}(x,y)}
\right\}.$$
Ahora, usando $J'_{\alpha+2n+2}(t) = {{\alpha+2n+2} \over
t}J_{\alpha+2n+2}(t)-J_{\alpha+2n+3}(t)$ obtenemos
$$\phi _{\alpha+2n+2} (x,y) = {{xy} \over {x^2-y^2}}\left\{ {J_{\alpha+2n+2}
(x)yJ'_{\alpha+2n+2} (y)-J_{\alpha+2n+2} (y)xJ'_{\alpha+2n+2} (x)} \right\},$$
y de aqu\'i se sigue el lema.

\bigskip

Esto nos permite expresar las sumas parciales de la serie de Fourier-Neumann
mediante
$$S_n(f,x) = T_1(f,x)-T_2(f,x)+T_{3,n}(f,x)-T_{4,n}(f,x)$$
siendo
$$T_1(f,x) = \int_0^\infty  {x^{1/2}y^{1/2}{{xJ_{\alpha +1}(x)J_\alpha (y)} \over
{x^2-y^2}}f(y)\,dy},$$
$$T_2(f,x) = \int_0^\infty  {x^{1/2}y^{1/2}{{yJ_\alpha (x)J_{\alpha +1}(y)} \over
{x^2-y^2}}f(y)\,dy},$$
$$T_{3,n}(f,x) = \int_0^\infty  {x^{1/2}y^{1/2}{{xJ'_\nu (x)J_\nu (y)}
\over {x^2-y^2}}f(y)\,dy},$$
$$T_{4,n}(f,x) = \int_0^\infty  {x^{1/2}y^{1/2}{{yJ_\nu (x)J'_\nu (y)}
\over {x^2-y^2}}f(y)\,dy}$$
y donde hemos denotado $\nu=\alpha+2n+2$.

De aqu\'i que para demostrar (2) bastar\'a con que probemos
$$\Vert T_if \Vert _p \le C \Vert f \Vert _p,\qquad i = 1,2 \eqno(8)$$
y
$$\Vert T_{i,n}f \Vert _p \le C \Vert f \Vert _p,\qquad i = 3,4. \eqno(9)$$

Es de destacar que los operadores $T_{3,n}f$ y $T_{4,n}f$ son muy parecidos a los
que aparecen en [1] y [2], aunque el m\'etodo que se utiliza aqu\'i permite
estudiar su acotaci\'on uniforme de modo m\'as sencillo. Siguiendo nuestro
m\'etodo se pueden simplificar f\'acilmente las demostraciones de [1] y [2].

En nuestro estudio emplearemos las siguientes acotaciones para funciones de
Bessel [7]
$$\alpha \ge -1/2 \quad \Longrightarrow \quad |J_\alpha (x)|\le
Cx^{-1/2},\quad |J_{\alpha +1}(x)|\le Cx^{-1/2}, \eqno(10)$$
$$-1 < \alpha < -1/2 \quad \Longrightarrow \quad |J_\alpha (x)|\le
Cx^{-1/2}(1+x^{\alpha +1/2}),\quad |J_{\alpha +1}(x)|\le C{x^{-1/2}
\over 1+x^{\alpha +1/2}}. \eqno(11)$$
Adem\'as,
$$|J_\nu (x)| \le Cx^{-1/4}\left( {|x-\nu |+\nu ^{1/3}} \right)^{-1/4},\quad
|J'_\nu (x)| \le Cx^{-3/4}\left( {|x-\nu |+\nu ^{1/3}} \right)^{1/4} \eqno(12)$$
donde, de nuevo, $\nu = \alpha+2n+2$. En todas ellas, la constante $C$
s\'olo puede depender de $\alpha$. Estas \'ultimas acotaciones son f\'aciles de
deducir de las que se utilizan en [1] y [2] y su expresi\'on final resulta
m\'as manejable.

Como \'ultima herramienta, necesitamos introducir la teor\'ia $A_p$ de pesos
(ver [3]). Dado $p \in (1,\infty)$, decimos que un peso $w$ definido en un
intervalo $(a,b)$ est\'a en $A_p(a,b)$ si
$$\left(\int_I w(x)\,dx\right)\left(\int_I w(x)^{-1/(p-1)}dx\,\right)^{p-1} \le
C|I|^p$$
con $C$ independiente de $I$ intervalo, $I\subseteq (a,b)$.

La importancia de la teor\'ia $A_p$ radica en su relaci\'on con la acotaci\'on
de la transformada de Hilbert
$$H(f,x) = V.P.\int_a^b {{{f(y)} \over {x-y}}\,dy}$$
(generalmente, omitimos $V.P.$ en su notaci\'on). En efecto, en [4] se
demuestra que
$$H:L^p((a,b),w) \longrightarrow L^p((a,b),w) {\rm\ acotada} \Longleftrightarrow
w \in A_p(a,b).$$
Adem\'as, la norma como operador de la transformada de Hilbert y la constante de
la definici\'on de $A_p$ dependen \'unicamente una de la otra, lo cual nos va a
permitir utilizar la teor\'ia $A_p$ uniforme.

Supongamos que una familia de pesos $\{w_n\}_{n=0}^\infty$ definidos en un
mismo intervalo $(a,b)$ cumplen la condici\'on $A_p$ con una misma constante $C$
(en este caso diremos que $w_n \in A_p(a,b)$ uniformemente). Entonces la familia
de transformadas de Hilbert $H:L^p((a,b),w_n) \longrightarrow L^p((a,b),w_n)$ es
uniformemente acotada.

Es bien conocido y sencillo de probar directamente el siguiente resultado
$$x^\beta \in A_p(0,1)\Longleftrightarrow x^\beta \in A_p(1,\infty
)\Longleftrightarrow -1<\beta <p-1. \eqno(13)$$
Adem\'as, mediante el cambio de variable $x=r_nz$ se demuestra f\'acilmente que
cuando $r_n\searrow 0$ entonces
$$\left( {|x|+r_n} \right)^\beta \in A_p(-1,1) {\rm\ uniformemente}
\Longleftrightarrow -1<\beta <p-1. \eqno(14)$$

Veamos a continuaci\'on varios resultados que utilizaremos m\'as adelante:

\bigskip

\proclaim LEMA 2.
\item{i)}$\displaystyle{-1<\alpha \le {-1\over 2},\ p>{2 \over {3+2\alpha
}}\Longrightarrow x^{p/2-1/2}\left( {1+x^{\alpha /2+1/4}} \right)^{-p}\in
A_p(0,\infty )}$
\item{ii)}$\displaystyle{-1<\alpha \le {-1\over 2},\ p<{2 \over {-1-2\alpha
}}\Longrightarrow x^{-1/2}\left( {1+x^{\alpha /2+1/4}} \right)^p\in
A_p(0,\infty )}$
\item{iii)}$\displaystyle{\alpha >-1,\;p>{4\over 3},\;0<s_n\nearrow \infty
\Longrightarrow x^{3p/8-1/2}\left( {|x^{1/2}-s_n|+s_n^{1/3}} \right)^{p/4}\in
A_p(0,\infty )}$ unif.
\item{iv)}$\displaystyle{\alpha >-1,\ p<4,\ 0<s_n\nearrow \infty \Longrightarrow
x^{p/8-1/2}\left( {|x^{1/2}-s_n|+s_n^{1/3}} \right)^{-p/4}\in A_p(0,\infty )}$
unif.

{\it Demostraci\'on}. i) y iii) son an\'alogos a ii) y iv) respectivamente,
luego vamos a probar \'unicamente estos dos \'ultimos.

\noindent ii) Si llamamos $w(x) = x^{-1/2}\left( {1+x^{\alpha /2+1/4}}
\right)^p$, no es dif\'icil demostrar que basta con comprobar $w \in A_p(0,1)$ y
$w \in A_p(1,\infty)$. Adem\'as, es claro que $w(x)$ se comporta como $x^{-1/2}$
en $(0,1)$ y como $x^{-1/2+p/4+\alpha p/4}$ en $(1,\infty)$ con lo cual, usando
(13) se sigue el resultado.

\noindent iv) Haciendo el cambio de variable $x=s_n^2z$ en la definici\'on de
$A_p$, es f\'acil ver que iv) es equivalente a demostrar
$$z^{p/8-1/2}\left( {|z^{1/2}-1|+s_n^{-2/3}} \right)^{-p/4}\in
A_p(0,\infty ) {\rm\ unif.}$$
De acuerdo al comportamiento de esta expresi\'on en los intervalos $(0,1/2)$,
$(1/2,3/2)$ y $(3/2,\infty)$, basta con que probemos
$$z^{p/8-1/2}\in A_p(0,{\textstyle{1 \over 2}}),\quad \left(
{|z^{1/2}-1|+s_n^{-2/3}} \right)^{-p/4}\in A_p({\textstyle{1 \over
2}},{\textstyle{3 \over 2}}),\quad z^{-1/2}\in A_p({\textstyle{3 \over
2}},\infty).$$
El primero y el tercero son cierto $\forall p\in(1,\infty)$ por (13), y el
segundo se deduce para $p<4$ usando $|z^{1/2}-1|\sim|z-1|$ y haciendo un cambio
de variable en (14).

\bigskip

Con todo esto, ya estamos en condiciones de abordar nuestro teorema principal

\bigskip

\proclaim TEOREMA.
Sean $\alpha>-1$, $1<p<\infty$ y $S_nf$ las sumas parciales de la serie de
Fourier-Neumann de $f$. Entonces, la acotaci\'on uniforme
$$\Vert S_nf \Vert _p \le C \Vert f \Vert _p,\qquad f \in L^p$$
es equivalente a que se verifiquen las desigualdades
$$4/3 < p < 4\quad {\sl cuando }\quad \alpha \ge -3/4$$
o
$${2 \over {3+2\alpha}} < p < {2 \over {-1-2\alpha}}\quad 
{\sl cuando } \quad \alpha < -3/4.$$

{\it Demostraci\'on}. De (5), (6) y (7) se sigue que las desigualdades sobre
$\alpha$ y $p$ son necesarias para la acotaci\'on uniforme. Veamos que
tambi\'en son suficientes, es decir, que se verifican (8) y (9). Analizaremos
\'unicamente $T_1f$ y $T_{4,n}f$ ya que $T_2f$ y $T_{3,n}f$ son an\'alogos.
En toda la demostraci\'on emplearemos siempre que sea necesario los cambios de
variable $x_1=x^2$ y $y_1=y^2$.
Para estudiar $T_1f$ hay que distinguir los casos $\alpha \ge -1/2$ y $\alpha <
-1/2$.

\noindent Empecemos con el caso $\alpha \ge -1/2$. Si denotamos
$g(y_1)=f(y_1^{1/2})y_1^{-1/4}J_\alpha(y_1^{1/2})$ y empleamos sucesivamente la
segunda acotaci\'on de (10), $x^{p/2-1/2} \in A_p(0,\infty)$ y la primera
acotaci\'on de (10) tenemos

$$\Vert T_1f \Vert^p_p = \int_0^\infty  {\left| {\int_0^\infty 
{x^{1/2}y^{1/2}{{xJ_{\alpha +1}(x)J_\alpha (y)} \over {x^2-y^2}}f(y)\,dy}}
\right|^p\,dx}$$

$$=2^{-p-1}\int_0^\infty  {\left| {H(g,x_1)} \right|^p\left| {J_{\alpha
+1}(x_1^{1/2})} \right|^px_1^{3p/4-1/2}\,dx_1}$$

$$\le C \int_0^\infty  {\left| {H(g,x_1)} \right|^p x_1^{p/2-1/2}\,dx_1} \le C_1
\int_0^\infty  {\left| {g(x_1)} \right|^p x_1^{p/2-1/2}\,dx_1}$$

$$\le C_2 \int_0^\infty  {\left| {f(x_1^{1/2})} \right|^p x_1^{-1/2}\,dx_1} =
C_3 \Vert f \Vert^p_p. $$

\noindent Para el caso $\alpha \ge -1/2$, con la misma notaci\'on para $g(y_1)$
y empleando sucesivamente la segunda acotaci\'on de (11), i) del lema 2 y la
primera acotaci\'on de (11) tenemos

$$\Vert T_1f \Vert^p_p = 2^{-p-1}\int_0^\infty  {\left| {H(g,x_1)}
\right|^p\left| {J_{\alpha +1}(x_1^{1/2})} \right|^px_1^{3p/4-1/2}\,dx_1}$$

$$\le C \int_0^\infty  {\left| {H(g,x_1)} \right|^p \left( {1+x_1^{\alpha
/2+1/4}} \right)^{-p} x_1^{p/2-1/2}\,dx_1} $$

$$\le C_1 \int_0^\infty  {\left|
{g(x_1)} \right|^p \left( {1+x_1^{\alpha /2+1/4}} \right)^{-p}
x_1^{p/2-1/2}\,dx_1}$$

$$\le C_2 \int_0^\infty  {\left| {f(x_1^{1/2})} \right|^p x_1^{-1/2}\,dx_1} =
C_3 \Vert f \Vert^p_p. $$

\noindent Por \'ultimo, estudiemos la acotaci\'on uniforme de $T_{4,n}f$.
Denotando $\nu=\alpha+2n+2$, $g(y_1)=f(y_1^{1/2})y_1^{1/4}J'_\nu(y_1^{1/2})$ y
utilizando sucesivamente la cota para $J_\nu$ de (12), iv) del lema 2 y la cota
para $J'_\nu$ de (12) obtenemos

$$\Vert T_{4,n}f \Vert^p_p = \int_0^\infty  {\left| {\int_0^\infty 
{x^{1/2}y^{1/2}{{yJ_\nu(x)J'_\nu(y)} \over {x^2-y^2}}f(y)\,dy}}
\right|^p\,dx}$$

$$=2^{-p-1}\int_0^\infty  {\left| {H(g,x_1)} \right|^p\left| {J_\nu(x_1^{1/2})}
\right|^px_1^{p/4-1/2}\,dx_1}$$

$$\le C \int_0^\infty  {\left| {H(g,x_1)} \right|^p x_1^{p/8-1/2}\left(
{|x_1^{1/2}-\nu|+\nu^{1/3}} \right)^{-p/4}\,dx_1} $$

$$\le C_1 \int_0^\infty  {\left| {g(x_1)} \right|^p x_1^{p/8-1/2}\left(
{|x_1^{1/2}-\nu|+\nu^{1/3}} \right)^{-p/4}\,dx_1}$$

$$\le C_2 \int_0^\infty  {\left| {f(x_1^{1/2})} \right|^p x_1^{-1/2}\,dx_1} =
C_3 \Vert f \Vert^p_p, $$
con lo que queda demostrado el teorema.


\bigskip

\centerline {\bf BIBLIOGRAF\'IA}

\bigskip
\frenchspacing

\item {{\bf [1]}} J. A. Barcel\'o y A. C\'ordoba, {\it Band-limited
functions: $L^p$-convergence}, Trans. Amer. Math. Soc. {\bf 313} (1989),
655--669.

\item {{\bf [2]}} A. C\'ordoba, {\it The disc multiplier}, Duke Math. Jour. {\bf
58} (1989), 21--27.

\item {{\bf [3]}} J. Garc\'\i a Cuerva y J. L. Rubio de Francia,
{\it ``Weighted Norm Inequalities and Related Topics''}, North-Holland,
Amsterdam, 1985.

\item {{\bf [4]}} R. Hunt, B. Muckenhoupt y R. Wheeden, {\it Weighted norm
inequalities for the conjugate function and Hilbert transform}, Trans. Amer.
Math. Soc. {\bf 176} (1973), 227--251.

\item {{\bf [5]}} J. Newman y W. Rudin, {\it Mean convergence of orthogonal
series}, Proc. Amer. Math. Soc. {\bf 3} (1952), 219--222.

\item {{\bf [6]}} H. Pollard, {\it The mean convergence of orthogonal
series I}, Trans. Amer. Math. Soc. {\bf 62} (1947), 387--403.

\item {{\bf [7]}} G. N. Watson, {\it ``A Treatise on the Theory of Bessel
Functions''}, Cambridge Univ. Press, 1966.

\item {{\bf [8]}} J. E. Wilkins, Jr., {\it Neumann series of Bessel functions},
Trans. Amer. Math. Soc. {\bf 64} (1948), 359--385.

\item {{\bf [9]}} J. E. Wilkins, Jr., {\it Neumann series of Bessel functions.
II}, Trans. Amer. Math. Soc. {\bf 69} (1950), 55--65.

\bigskip

\noindent{\sl Juan Luis Varona Malumbres}\par
\noindent{\sl Departamento de Matem\'atica Aplicada}\par
\noindent{\sl Colegio Universitario de La Rioja (Universidad de Zaragoza)}\par
\noindent{\sl Obispo Bustamante, 3}\par
\noindent{\sl 26001 Logro\~no, ESPA\~NA}


\end

